\section{From loss to decoding: the logit-margin lemma}
\label{sec:logit_margin_decoding}

\begin{lemma}[Logit-margin sufficient condition for correct decoding]\label{lem:logit_margin_decoding}
Fix $Q \in F$ and $x \in \R^d$. The following two implications hold.
\begin{enumerate}[label=(\roman*),leftmargin=*]
\item If $\loss(x; Q) \le \log 2$ then $\cmass(x; Q) \ge 1/2$.
\item If, in addition, the non-degeneracy condition
\begin{align}\label{eq:logit_margin_nondegen}
   |\Cset(Q)| \cdot \max_{v \notin \Cset(Q)} p(x)_v
   \;<\;
   \max_{v \in \Cset(Q)} p(x)_v,
\end{align}
holds, then $\dec(x) \in \Cset(Q)$.
\end{enumerate}
\end{lemma}

\begin{proof}
The two parts are proved separately.

\noindent\textbf{Part (i).}
By \Cref{eq:loss_def}, $\loss(x; Q) = -\log \cmass(x; Q)$. Therefore
$\loss(x; Q) \le \log 2$ is equivalent to $-\log \cmass(x; Q) \le \log 2$,
which, applying the (strictly decreasing) function $t \mapsto e^{-t}$
to both sides and rearranging, is equivalent to
$\cmass(x; Q) \ge e^{-\log 2} = 1/2$. This is the assertion of (i).

\noindent\textbf{Part (ii).}
Assume both $\cmass(x; Q) \ge 1/2$ (from (i)) and the non-degeneracy
condition \Cref{eq:logit_margin_nondegen}. We must show
$\dec(x) = \arg\max_v (W_U x)_v \in \Cset(Q)$. Since
$p(x) = \softmax(W_U x)$ is a strictly monotone transformation of
$W_U x$ (the softmax preserves the order of its inputs), the argmax over
$v$ of $(W_U x)_v$ coincides with the argmax of $p(x)_v$. It therefore
suffices to show
\begin{align}\label{eq:logit_margin_p_argmax}
   \max_{v \in \Cset(Q)} p(x)_v
   \;>\;
   \max_{v \notin \Cset(Q)} p(x)_v.
\end{align}
This is implied directly by \Cref{eq:logit_margin_nondegen}:
$|\Cset(Q)| \cdot \max_{v \notin \Cset(Q)} p(x)_v < \max_{v \in \Cset(Q)} p(x)_v$
and $|\Cset(Q)| \ge 1$ together yield
$\max_{v \notin \Cset(Q)} p(x)_v < \max_{v \in \Cset(Q)} p(x)_v$, which is
exactly \Cref{eq:logit_margin_p_argmax}. Hence $\dec(x) \in \Cset(Q)$.
\end{proof}

\begin{remark}\label{rem:logit_margin_nondegen_typical}
The non-degeneracy condition \Cref{eq:logit_margin_nondegen} is a mild
``no two-way tie'' condition between the correct and incorrect tokens.
A sufficient one-line version: when $\cmass(x; Q) \ge 1/2$, by
pigeonhole $\max_{v \in \Cset(Q)} p(x)_v \ge \cmass(x; Q)/|\Cset(Q)| \ge 1/(2 |\Cset(Q)|)$,
so \Cref{eq:logit_margin_nondegen} is implied by
$\max_{v \notin \Cset(Q)} p(x)_v < 1/(2 |\Cset(Q)|^2)$ --- a strong
``no incorrect token captures non-trivial mass'' condition. For typical
trained reasoning policies on verifiable tasks the incorrect-token
posteriors are exponentially small in the logit margin, so the
condition is satisfied wherever $\loss(x; Q) \le \log 2$.
\end{remark}

\begin{remark}[Sufficient, not exact: the polyhedral cone is larger]\label{rem:sufficient_not_exact}
The level set $\{x : \loss(x; Q) \le \log 2\}$ is a smooth, convex-by-level
subset of $\R^d$, while the exact decoding-correctness region is the
polyhedral cone
\begin{align*}
   \mathcal D(Q) \;\coloneqq\; \bigl\{\, x \in \R^d :
   \max_{v \in \Cset(Q)} (W_U x)_v
   \;\ge\;
   \max_{v \notin \Cset(Q)} (W_U x)_v \,\bigr\},
\end{align*}
the (closed, polyhedral) set on which the argmax over $\mathcal V$ falls
within $\Cset(Q)$. \Cref{lem:logit_margin_decoding} establishes a smooth
sufficient condition for $x \in \mathcal D(Q)$ via the logit-margin
threshold $\loss(x; Q) \le \log 2$; that is, the level set
$\{\loss \le \log 2\}$ is a subset of $\mathcal D(Q)$, but \emph{not}
all of $\mathcal D(Q)$. The cone $\mathcal D(Q)$ extends far outside
the level set: for any $x$ whose logits already place the argmax in
$\Cset(Q)$ by an arbitrarily small margin, $x \in \mathcal D(Q)$ even
when $\loss(x; Q) \gg \log 2$. A tighter analysis would replace the
smooth-loss surrogate by a direct treatment of the cone $\mathcal D(Q)$
(e.g.\ via a piecewise-linear margin function). We retain the smooth
surrogate because it is what couples cleanly to the descent analysis of
\Cref{lem:descent_inequality,lem:telescoping}; the cone-based treatment
is a natural future direction.
\end{remark}
